%% Chapter 6 — Section 6.5: Exercises
\section{Exercises}
\label{sec:6.5}

\begin{enumerate}

\item[\textbf{6.1}]
Consider the quadratic objective function $V(x) = \frac{1}{2}(x-\mu)^T H(x-\mu)$ with
positive definite $H$. Denote by $\lambda_{\min}(H)$ and $\lambda_{\max}(H)$ the smallest and
largest eigenvalues of $H$. Show that the gradient descent algorithm given by
\[
  x_{n+1} = x_n - \epsilon\nabla V(x_n), \qquad n = 0, 1, \ldots
\]
with step-size $\epsilon < 2\bigl(\lambda_{\min}(H) + \lambda_{\max}(H)\bigr)^{-1}$ converges
linearly to the minimizer $\mu$. Specifically, show that
\[
  \norm{x_n - \mu} \le \bigl(1 - \epsilon\lambda_{\min}(H)\bigr)^n \norm{x_0 - \mu}.
\]

\item[\textbf{6.2}]
In this exercise, you will compare MALA and RWMH algorithms by sampling from Gaussian
distributions, focusing on how tuning parameters affect their performance, acceptance rate, and
state space exploration.

\begin{enumerate}[label=\alph*)]
  \item \textbf{MALA with a Standard Gaussian Target.}
    Generate $10{,}000$ samples from a standard normal distribution, $f = \Normal(0,1)$, using
    MALA. Investigate the impact of varying the step-size $\epsilon$ on the algorithm's
    performance. Complete the following tasks:
    \begin{itemize}
      \item Find a step-size $\epsilon$ that yields an acceptance rate close to $0.574$.
      \item Create a trace plot for the chain with the step-size you found in the previous
        step, along with trace plots with $\epsilon = 0.001$ and $\epsilon = 10$. These plots
        should illustrate different exploration dynamics.
      \item For the selected $\epsilon$, compare the empirical sample distribution with the
        standard normal distribution through histograms.
      \item Discuss how $\epsilon$ influences exploration efficiency and acceptance rate
        trade-offs.
    \end{itemize}

  \item \textbf{RWMH with a Standard Gaussian Target.}
    Using the same setup as in part a), explore RWMH performance by varying the standard
    deviation parameter, $\sigma$ in the proposal distribution $g_\sigma = \Normal(0, \sigma^2)$.
    Repeat the tasks described in part a), substituting step-size $\epsilon$ with $\sigma$ and
    aiming to plot results similar to those in Figure~4.1.

  \item \textbf{High-Dimensional Gaussian Target.}
    Consider a 10-dimensional Gaussian distribution $f = \Normal(0, I_{10})$, where $0$ is a
    10-dimensional zero vector and $I_{10}$ is the identity matrix of size $10$. Find a choice
    of $\epsilon$ to achieve an acceptance rate close to $0.574$, and find the $\sigma$ for
    RWMH that matches its optimal higher-dimensional acceptance probability (about $0.234$).
    With these parameters, plot the mean estimates from MALA and RWMH for all marginals of
    this 10-dimensional Gaussian, comparing them against the true mean in a histogram.
\end{enumerate}

\item[\textbf{6.3}]
We say that a response random variable $Y$ given explanatory variables $(X_1, \ldots, X_d)$
follows a logistic regression model if
\[
  Y = \begin{cases} 1 & \text{with probability } p, \\ 0 & \text{with probability } 1-p, \end{cases}
\]
where
\[
  p = \frac{1}{1 + \exp\bigl(-(w_0 + w_1 X_1 + \ldots + w_d X_d)\bigr)}.
\]
In this exercise, you will use MALA to estimate the vector $w = (w_0, \ldots, w_d)$ using $K$
independent pairs $\{(X^{(k)}, Y^{(k)})\}_{k=1}^K$ of response and explanatory variables.
Notice that $X^{(k)} = (X_1^{(k)}, \ldots, X_d^{(k)}) \in \R^d$ denotes the $k$-th
explanatory variable and $Y^{(k)} \in \{0,1\}$ denotes the $k$-th response variable.

\begin{enumerate}[label=\arabic*.]
  \item \textbf{Setting and Preparation:} Set $d = 2$ and $w = (0, 5, -5)$. Simulate $K = 500$
    explanatory variables $\{X^{(k)}\}_{k=1}^K = \{(X_1^{(k)}, X_2^{(k)})\}_{k=1}^K$ by
    uniformly sampling in $[0,1] \times [0,1]$. For each $X^{(k)}$, simulate a label $Y^{(k)}$
    according to the logistic regression model.

    \begin{enumerate}[label=\alph*)]
      \item Plot the response variables $\{(X^{(k)})\}_{k=1}^K$, marking in red all points
        with label 1 and in black those with label 0. Describe the plot you generate.
    \end{enumerate}

  \item \textbf{Metropolis-Adjusted Langevin Algorithm:} Adopt a Bayesian formulation,
    specifying a Gaussian prior $w \sim \Normal(0, \sigma^2 I)$. You will implement MALA to
    sample from the posterior distribution of $w$ given the data
    $\mathbf{X} := \{X^{(k)}\}_{k=1}^K$ and $\mathbf{Y} := \{Y^{(k)}\}_{k=1}^K$.

    \begin{enumerate}[label=\alph*),start=2]
      \item Show that the log-likelihood is given by
        \begin{align*}
          \log f(\mathbf{Y} \mid \mathbf{X}, w)
          &= -\sum_{k=1}^K Y^{(k)} \log\!\bigl(1 + \exp(-Z^{(k)})\bigr) \\
          &\quad - \sum_{k=1}^K (1-Y^{(k)})\Bigl[Z^{(k)} + \log\!\bigl(1 + \exp(-Z^{(k)})\bigr)\Bigr],
        \end{align*}
        where $Z^{(k)} = w_0 + w_1 X_1^{(k)} + \ldots + w_d X_d^{(k)}$.

      \item Show that the log-posterior density is given by
        \begin{align*}
          \log f(w \mid \mathbf{X}, \mathbf{Y})
          &= -\sum_{k=1}^K Y^{(k)} \log\!\bigl(1 + \exp(-Z^{(k)})\bigr) \\
          &\quad - \sum_{k=1}^K (1-Y^{(k)})\Bigl[Z^{(k)} + \log\!\bigl(1 + \exp(-Z^{(k)})\bigr)\Bigr] \\
          &\quad - \frac{1}{2\sigma^2}\sum_{j=0}^d w_j^2 + c,
        \end{align*}
        where $c$ is a constant.

      \item Compute $\nabla_w \log f(w \mid \mathbf{X}, \mathbf{Y})$.

      \item Show that the log-acceptance probability with which MALA accepts a move from
        $w^{(n)}$ to a proposed value $w^*$ is given by
        \begin{align*}
          \log\!\bigl(a(w^{(n)}, w^*)\bigr)
          = \min\Bigl\{0,\;
          &\log f(w^* \mid \mathbf{X}, \mathbf{Y}) - \log f(w^{(n)} \mid \mathbf{X}, \mathbf{Y}) \\
          &+ \frac{\norm{w^* - w^{(n)} - \epsilon\nabla_w \log f(w^{(n)} \mid \mathbf{X}, \mathbf{Y})}^2}{4\epsilon} \\
          &- \frac{\norm{w^{(n)} - w^* - \epsilon\nabla_w \log f(w^* \mid \mathbf{X}, \mathbf{Y})}^2}{4\epsilon}
          \Bigr\}.
        \end{align*}

      \item Based on the above derivations, implement MALA to sample from the posterior
        distribution of $w$. Set the prior variance parameter to $\sigma^2 = 1$. Tune the
        step-size to achieve an acceptance rate close to $0.574$, and document the final
        step-size selected. To validate your implementation and showcase the convergence of
        the algorithm, plot the trajectories of MALA for $w_0$, $w_1$, and $w_2$, comparing
        them with the true coefficients $(0, 5, -5)$.
    \end{enumerate}
\end{enumerate}

\item[\textbf{6.4}]

\begin{enumerate}[label=\alph*)]
  \item Replicate the results reported in Figure~\ref{fig:6.3} by simulating the trajectories
    of $10{,}000$ particles under Langevin dynamics. Provide plots of your results.

  \item For $\beta > 0$, we can generalize Langevin equation \eqref{eq:6.7} into
    \begin{equation}
      \label{eq:6.9}
      dX(t) = -\nabla V\!\bigl(X(t)\bigr)\,dt + \sqrt{2\beta^{-1}}\,dW, \quad 0 \le t < \infty.
    \end{equation}
    As will be discussed in Chapter~7, the parameter $\beta$ can then be interpreted as an
    inverse temperature. Show that the process defined by the generalized Langevin equation
    \eqref{eq:6.9} has $f_\beta(x) \propto \exp(-\beta V(x))$ as invariant distribution.
    For the double-well potential $V(x) = x^4 - x^2 + 0.25$, complete the following tasks:
    \begin{itemize}
      \item On a single plot, graph the scaled potential functions $\beta V(x)$ for
        $\beta = 1, 10, 50$.
      \item For each choice of $\beta$, identify a step-size $\epsilon$ that yields an
        acceptance rate close to $0.574$.
      \item For each $\beta$ and corresponding step-size $\epsilon$, repeat the MALA simulation
        $100$ times starting near $x = -\frac{1}{\sqrt{2}}$ to determine the average steps
        taken to first approach $x = 0$ and examine the algorithm's ability to explore the
        state space.
      \item Discuss how the value of $\beta$ affects particle transitions between the two wells
        of the potential.
    \end{itemize}
\end{enumerate}

\end{enumerate}
